\documentclass[11pt]{article}
\usepackage[margin=1in]{geometry}
\usepackage{amsmath,amssymb,amsthm}
\usepackage{enumitem}

\newtheorem{theorem}{Theorem}[section]
\newtheorem{lemma}[theorem]{Lemma}
\newtheorem{proposition}[theorem]{Proposition}
\newtheorem{remark}[theorem]{Remark}

\newcommand{\Om}{\Omega}
\newcommand{\R}{\mathbb R}
\newcommand{\abs}[1]{\left|#1\right|}
\newcommand{\dH}{\,d\mathcal H^{N-1}}

\title{Selection as Calibrated Symmetrization}
\author{}
\date{May 2026}

\begin{document}
\maketitle

\section{Point of the note}

The selection principle in Plan 1 looks like a penalized minimization, but it
has a useful geometric interpretation.  It is an implicit, constrained
symmetrization step:

\[
   \hbox{replace } \Omega_0
   \hbox{ by the best torsion-energy representative with essentially the same }
   \alpha\hbox{-asymmetry.}
\]

It is not Schwarz or Steiner symmetrization.  Those flows drive the asymmetry
to zero, while the selection principle must preserve enough asymmetry to keep
a near-counterexample alive.  The correct object is therefore a
\emph{calibrated symmetrization}: decrease the torsion deficit and volume
error, but constrain the smooth asymmetry level.

This point of view also explains why the perturbing functional can be taken to
be Lipschitz, with uniformly bounded boundary Lipschitz constant.

\section{The selected resolvent}

Let \(\Omega_0\subset B_R\), \(\abs{\Omega_0}=\abs{B_1}\), and set
\[
   \varepsilon:=\alpha(\Omega_0),
   \qquad
   \delta:=E(\Omega_0)-E(B_1),
   \qquad
   q:=\delta/\varepsilon.
\]
Recall the BDV smooth asymmetry
\[
   \alpha(U)
      =
      \int_{U\Delta B_1(x_U)}
        \bigl|1-\abs{x-x_U}\bigr|\,dx,
\]
where \(x_U\) is the barycenter.  On bounded sets,
\[
   \abs{\alpha(U)-\alpha(V)}
      \le C_R\abs{U\Delta V}.
   \tag{2.1}\label{eq:alpha-l1-lip}
\]

For a parameter \(\tau\ll1\), define
\[
   \Phi_{\varepsilon,\tau}(s)
      :=
      \sqrt{\varepsilon^2+\tau^2(s-\varepsilon)^2}.
\]
Then
\[
   0\le \Phi_{\varepsilon,\tau}'(s)\le \tau
   \quad\hbox{in absolute value.}
   \tag{2.2}\label{eq:phi-slope}
\]

The selected set is
\[
   S_{\tau}(\Omega_0)
      \in
      \operatorname*{argmin}_{U\subset B_R}
      \left\{
      E(U)+f_{\hat\eta}(\abs U)
      +\Phi_{\varepsilon,\tau}(\alpha(U))
      \right\}.
   \tag{2.3}
\]
This is exactly the Plan 1/BDV penalized minimizer, with the sequence parameter
replaced by the single parameter \(\tau\).

\begin{proposition}[Calibrated symmetrization properties]\label{prop:cal}
Let \(U=S_\tau(\Omega_0)\), and let
\[
   F_{\hat\eta}(U):=E(U)+f_{\hat\eta}(\abs U).
\]
Then
\[
   F_{\hat\eta}(U)\le F_{\hat\eta}(\Omega_0).
   \tag{2.4}\label{eq:energy-improve}
\]
Moreover
\[
   \frac{\abs{\alpha(U)-\varepsilon}}{\varepsilon}
      \le
      \frac{\sqrt{2q+q^2}}{\tau}
      +O_R(q),
   \tag{2.5}\label{eq:alpha-preserve}
\]
with the harmless \(O_R(q)\) coming from the usual volume-penalty
normalization.
\end{proposition}

\begin{proof}
Minimality gives
\[
   F_{\hat\eta}(U)+\Phi_{\varepsilon,\tau}(\alpha(U))
      \le
   F_{\hat\eta}(\Omega_0)+\Phi_{\varepsilon,\tau}(\varepsilon)
      =
   F_{\hat\eta}(\Omega_0)+\varepsilon.
\]
Since \(\Phi_{\varepsilon,\tau}\ge\varepsilon\), this implies
\eqref{eq:energy-improve}.  Also, using the lower bound
\(F_{\hat\eta}(U)\ge F_{\hat\eta}(B_1)-O_R(\delta)\) supplied by the BDV
volume-penalty lemma, the same inequality gives
\[
   \sqrt{\varepsilon^2+\tau^2(\alpha(U)-\varepsilon)^2}
      \le
   \varepsilon+\delta+O_R(\delta\varepsilon).
\]
Squaring and dividing by \(\varepsilon^2\) gives
\eqref{eq:alpha-preserve}.
\end{proof}

Thus the map \(S_\tau\) really is an improvement step: it lowers the
volume-penalized torsion energy, while the soft constraint prevents the
asymmetry from collapsing.  With the canonical choice \(\tau=q^{1/4}\), the
relative asymmetry loss is \(O(q^{1/4})\).

\section{The flow behind the minimization}

For a smooth deformation with normal velocity \(V\), the torsion energy has
the shape derivative
\[
   E'(U)[V]
      =
      -\frac12\int_{\partial U}\abs{\nabla u_U}^2 V\,\dH.
   \tag{3.1}
\]
The volume term contributes \(f_{\hat\eta}'(\abs U)\int_{\partial U}V\,\dH\).
The smooth asymmetry has first variation
\[
   \alpha'(U)[V]
      =
      \int_{\partial U}\Psi_U V\,\dH,
   \tag{3.2}
\]
where, up to an additive constant absorbed into the volume multiplier,
\[
   A_U:=\frac1{\abs U}
      \int_U\frac{z-x_U}{\abs{z-x_U}}\,dz,
   \qquad
   \Psi_U(x):=\abs{x-x_U}-A_U\cdot(x-x_U).
   \tag{3.3}
\]
Since \(\abs{A_U}\le1\), the bracket is globally Lipschitz:
\[
   \operatorname{Lip}(\Psi_U)\le 2.
   \tag{3.4}
\]

At a smooth selected minimizer, stationarity gives the Bernoulli law
\[
   \abs{\nabla u_U(x)}^2
      =
      \Lambda
      +
      a\,\Psi_U(x),
   \qquad x\in\partial U,
   \tag{3.5}\label{eq:bernoulli-law}
\]
where
\[
   a=2\Phi_{\varepsilon,\tau}'(\alpha(U)),
   \qquad
   \abs a\le 2\tau,
   \tag{3.6}
\]
up to the harmless convention-dependent factor \(2\).  Hence the boundary
coefficient in \eqref{eq:bernoulli-law} has a Lipschitz constant bounded by
\(C\tau\).
Once the free boundary lies in a fixed annulus around the selected center,
\(\Psi_U\) has \(C^k\) bounds depending only on \(N,R,k\), because the only
nonlinear term is \(\abs{x-x_U}\).

This is the flow interpretation: the formal normal velocity
\[
   V
      =
      \frac12\abs{\nabla u_U}^2
      -f_{\hat\eta}'(\abs U)
      -\Phi_{\varepsilon,\tau}'(\alpha(U))\Psi_U
   \tag{3.7}
\]
is a constrained torsion-symmetrizing descent.  The selected minimizer is a
stationary point of this descent.  The direct minimization is an implicit
Euler/resolvent version of the flow, and it avoids having to construct the
flow itself.

\section{Local replacement version}

The same idea can be phrased without global minimization.  Given a current
positivity set \(U=\{u>0\}\), choose a ball \(B\subset B_R\), keep \(u\) fixed
outside \(B\), and replace \(u\) inside \(B\) by a competitor \(v\) that lowers
\[
   \frac12\int\abs{\nabla v}^2-\int v
      + f_{\hat\eta}(\abs{\{v>0\}})
      + \Phi_{\varepsilon,\tau}(\alpha(\{v>0\})).
\]
Because of \eqref{eq:alpha-l1-lip}--\eqref{eq:phi-slope}, the last term
changes by at most
\[
   C_R\tau\,
   \abs{\{u>0\}\Delta\{v>0\}}.
   \tag{4.1}
\]
Therefore any limit of a descent-by-local-replacements procedure satisfies the
one-phase quasi-minimality inequality
\[
\begin{aligned}
   \frac12\int\abs{\nabla u}^2-\int u
      + f_{\hat\eta}(\abs{\{u>0\}})
   \le{}&
   \frac12\int\abs{\nabla v}^2-\int v
      + f_{\hat\eta}(\abs{\{v>0\}}) \\
   &+ C_R\tau\,
      \abs{\{u>0\}\Delta\{v>0\}}.
\end{aligned}
\tag{4.2}
\]
This is precisely the Alt--Caffarelli/BDV input used for density,
nondegeneracy, Lipschitz control of \(u\), and eventual free-boundary
regularity.  In this local form the selection step looks like a smoothing
algorithm: keep making torsion-improving replacements, but charge any
asymmetry drift by a uniformly Lipschitz cost.

\section{What this does and does not replace}

\begin{enumerate}[label=\arabic*.]
\item It does replace the vague compactness language by a concrete
``symmetrizing resolvent'':
\[
   \Omega_0\mapsto S_\tau(\Omega_0).
\]
The map lowers volume-penalized torsion energy and keeps the smooth asymmetry
within a controlled factor.

\item It explains why the extra functional can be chosen with bounded
Lipschitz constant.  The scalar penalty has slope at most \(\tau\), and the
asymmetry bracket has globally bounded Lipschitz norm on \(B_R\).  On the
eventual annular collar the bracket has all \(C^k\) bounds needed for the
Schauder bootstrap.

\item It is not a literal Schwarz/Steiner symmetrization.  Such
symmetrizations improve the energy by moving toward the ball, but they also
erase the asymmetry.  Selection must improve the set while keeping it at the
same asymmetry level, because otherwise a near-counterexample disappears.

\item It does not by itself prove graph entry.  It gives a regular
quasi-minimizer with a controlled Bernoulli coefficient.  The later Plan 1
steps still need density, Hausdorff-to-flatness, and patching to a global
graph.
\end{enumerate}

\section{Practical takeaway}

The underlying structure of the Plan 1 selection step is:
\[
   \boxed{
   \hbox{torsion symmetrization constrained to the asymmetry level of the input}
   }
\]
implemented by an implicit variational step.  For the purpose of regularity,
the only property of the asymmetry constraint that matters is
\[
   \abs{\Phi_{\varepsilon,\tau}(\alpha(U))
       -\Phi_{\varepsilon,\tau}(\alpha(V))}
      \le C_R\tau\abs{U\Delta V},
\]
and, at a smooth boundary,
\[
   \abs{\nabla u_U}^2=\Lambda+a\Psi_U,
   \qquad
   \abs a\le C\tau,
   \qquad
   \operatorname{Lip}(\Psi_U)\le C.
\]
So yes: after selection, one may treat the added functional as a Lipschitz
boundary forcing with bounded Lipschitz constant.  The penalized minimization
is the rigorous surrogate for the desired smoothing/flow-like symmetrization.

\end{document}
